\documentclass[output=paper]{langscibook}
\ChapterDOI{10.5281/zenodo.22880280}
\title{\textit{Kave goroto} – {The} turtle net}
\author{Mark Tahi and Enoch Horai Magum}
\abstract{}
\begin{document}
\maketitle\label{chap:27}
\begin{figure}[H]
    \includegraphics[trim=0 0 0 0,clip,width=.9\linewidth]{figures/27_1.jpg}
    \caption{\textit{Kave goroto} -- Turtle net}
    \label{fig:27_1}
\end{figure}
\begin{xltabular}{\textwidth}{lXX}
	(\ref{ex:text20-1})&\textit{\textbf{Kave goroto}}&\textbf{Turtle net}\\
	(\ref{ex:text20-2})&\textit{O kave goroto vai Teapu na pakupaku raara bono naono to dao raara bono baneo.}&The Teop turtle net, we make it from a tree that we call the breadfruit tree.\\
	(\ref{ex:text20-3})&\textit{O baneo bona he toro sii vaamanae koa.}&But the breadfruit tree must be quite young.\\
	(\ref{ex:text20-4})&\textit{A taba momohu toro paku ara, eara toro mosi gunaha o baneo.}&The first thing we must do, we must cut down the breadfruit tree.\\
	(\ref{ex:text20-5})&\textit{Beve obete nana teo kasuana, eara repaa kahu a kapa nae bono kehaa to dao raara bono sui.}&When it is lying on the ground, we scrape its bark off with a shell that we call \textit{sui} `ark clam'.\\
	(\ref{ex:text20-6})&\textit{Beara kahu vakavakavara e, eara repaa varoava e bono buaku ge o kukan o bon.}&When we have finished scraping it, we put it into the sun for two or three days.\\
\end{xltabular}
\begin{figure}[H]
    \captionsetup[subfigure]{oneside,margin={0cm,0cm}}
    \subfloat[]{\includegraphics[trim=0 35 0 0,clip,width=.47\linewidth]{figures/27_2.jpg}\label{fig:27_2}}\hfil
    \subfloat[]{\includegraphics[trim=0 0 0 0,clip,width=.53\linewidth]{figures/27_5.jpg}\label{fig:27_5}}
    \caption{(a) \textit{Koriitoo} -- Tiger Cowrie, (b) \textit{Sui} -- Ark Clam}
\end{figure}
\begin{figure}[H]
    \includegraphics[trim=0 55 0 0,clip,width=.42\linewidth]{figures/27_3.jpg}
    \includegraphics[trim=0 0 0 0,clip,width=.57\linewidth]{figures/27_4.jpg}
    \caption{\textit{Roravita} -- Helmet Shell}
    \label{fig:24_3-4}
\end{figure}
\begin{xltabular}{\textwidth}{lXX}
	(\ref{ex:text20-7})&\textit{Beve raka vaantee, eara repaa hihi kahi e bono naono nae.}&When it is dry enough, we remove the fibres from its wood.\\
	(\ref{ex:text20-8})&\textit{Beara hihi vakavara e, eara repaa hupihupi e bono sui.}&When we have removed them, we comb them with the ark shell.\\
	(\ref{ex:text20-9})&\textit{Beara hupihupi vakavara e, a kiu vaa tea vinugu repaa taneo, ore tavus potee bona kanono.}&When we have finished combing them, the work of twisting starts, so that they become a string.\\
	(\ref{ex:text20-10})&\textit{A kiu vai na antee nana bona buaku ge a dorana a otei tea peha kanono,}&This work needs two or four men for one string,\\
	(\ref{ex:text20-11})&\textit{bea meha buaku ge a dorana vinugu batari bona meha kanono.}&while another two or four twist another string.\\
	(\ref{ex:text20-12})&\textit{A kiu vai toro vatotoena ni ara, tabae a kave toro bebeahu.}&This work we must share, because the net should become long.\\
	(\ref{ex:text20-13})&\textit{Bea kiu vaa tea vinugu kavara, eara repaa vavaantee nia maa kanono,}&When the work of twisting is finished, we have made the strings long enough,\\
	(\ref{ex:text20-14})&\textit{ei kou toro bebeahu vaantee.}&because they must be long enough.\\
	(\ref{ex:text20-15})&\textit{Bea maa kanono vavaantee vakavara, eara repaa vaahana riori teo naono,}&When the strings have all become long enough, we hang them on a tree,\\
	(\ref{ex:text20-16})&\textit{eara repaa rahi vatotobin riori.}&then we pull them straight.\\
	(\ref{ex:text20-17})&\textit{Beve kavara o hum kiu bona, eara repaa taneo mana tea huha.}&When this part of the work is finished, we eventually start with the tying.\\
	(\ref{ex:text20-18})&\textit{A kave na tanetaneo raara bona taba to dao raara bona tavara.}&The net, we start with a thing that we call \textit{tavara}.\footnote{The \textit{tavara} is a rectangular piece of \textit{kururu} bamboo that is used to measure the size of the meshes, see Chapter \ref{chap:29}:\ref{ex:text22-13}.}\\
	(\ref{ex:text20-19})&\textit{Evehee, beara mei rara tea taneo, eara toro varaarare ni roho a tatabonata na kave bara bebeahu nae.}&But before we start, we must decide on the width of the net and its length.\\
\end{xltabular}	
\begin{xltabular}{\textwidth}{lXX}
	(\ref{ex:text20-20})&\textit{Beara vuha vaapeha meo ipinoo to paa vamoromoroko ni vuara,}&When we have agreed on the measurement that we have talked about,\\
    (\ref{ex:text20-21})&\textit{eara paku vai a tavara to kahi ipinoo bata nom ara o mata kave.}&we make the \textit{tavara} (a piece of bamboo) which we are going to use for mesuring the meshes of the net.\\
	(\ref{ex:text20-22})&\textit{Be kavara amaa vamoromoroko, eara repaa taneo tea huha kave.}&When the discussions are finished, we start tying the net.\\
	(\ref{ex:text20-23})&\textit{A kiu a beera, me kahi huha vuara bona maa bon amaa bero.}&The work is big, and we will tie (the net) many days.\\
	(\ref{ex:text20-24})&\textit{Be kavara a kiu vaa tea huha, eara repaa paku pete a kiu vaa tea ibun bono kehaa.}&When the work of tying is finsihed, we will continue with the work of threading the shells.\\
	(\ref{ex:text20-25})&\textit{Amaa kehaa bona to dao raara bono koriitoo.}&These shells we call \textit{koriitoo} `Tiger Cowrie'\\
	(\ref{ex:text20-26})&\textit{O koriitoo kou to kahi ibun raara teo roveana vaa paana kave, eove to kahi rahirahi vaadovo bata ni nom bona kave.}&The Tiger Cowry shells we thread onto the seam at the bottom of the net, they will pull the net down and make it sink into the deep sea.\\
	(\ref{ex:text20-27})&\textit{Be kavara a kiu bona, eara repaa gono pete maa o buo roravita mene vaataena kave.}&When this work is finished, we get two helmet shells for the bottom of the net.\\
	(\ref{ex:text20-28})&\textit{O peho teo meho punete, o meho he teo meho punete.}&One for one edge, the other one for the other edge.\\
	(\ref{ex:text20-29})&\textit{A bua kehaa bona to kahi ato vamasuu bata ni nom bona kave}&These two shells will hold the net steady\\
	(\ref{ex:text20-30})&\textit{to mene orooro kahi batana bona maa hum to kahi vavaagunaha kave raara.}&that otherwise would float away from the places where we would drop the net down.\\
	(\ref{ex:text20-31})&\textit{O buo kara haraharaa mee kahi pam ni raara teo buo roravita.}&We tie two feeling lines to the Helmet Shells.\\
	(\ref{ex:text20-32})&\textit{O ogaoga na kave he na paku raara bono naono no baneo to paa hihi kahi vuara bona kapa nae.}&The floats of the net, however, we make from the wood of the breadfruit tree from which we have removed its bark.\\
	(\ref{ex:text20-33})&\textit{Na mosimosi vakakato raara beara ipinoo bata rae bona ta naono to paa vaantee vuara bona bebeahu no buo kubisi bebeahu na kuri raara.}&We cut it into short pieces, while measuring it with a piece of wood that we have made as long as the length of the two middle fingers of our hand.\\
	(\ref{ex:text20-34})&\textit{Beara mosi kurukurus bata rae, a mosi toro kerun tamuana batana.}&When we cut it into pieces, the cut must always be diagonal.\\
	(\ref{ex:text20-35})&\textit{Eara kahi mosi kurus bata nom, beo meho sii aba vaanana vaaraka bata ni rori bari teo suraa.}&We will do the cutting, while other people are warming and drying the pieces of wood over the fire.\\
	(\ref{ex:text20-36})&\textit{Be kavara, eara repaa ibun mana riori tea kara kanono to dao raara bona roveana.}&When this is finished, we thread them onto the string that we call \textit{roveana} `the seam'.\\
	(\ref{ex:text20-37})&\textit{O kanono bona he na pakupaku raara bona roge.}&This string we make from the coastal hibiscus.\footnote{See Chapter \ref{chap:12}:\ref{ex:text8-6}.}\\
\end{xltabular}
\end{document}